
\section*{5. Restricted \texorpdfstring{\(p\)}{p}-Lie Algebras}
\phantomsection
\label{sga3:VIIA:5}
\addcontentsline{toc}{section}{5. Restricted p-Lie Algebras}

We first recall some results from the \emph{S\'eminaire Sophus
Lie}.

\medskip
\noindent\textbf{5.1.}
\phantomsection
\label{sga3:VIIA:5.1}

Let \(p\) be a prime number, let \(R\) be a commutative ring of
characteristic \(p\), and let \(A\) be an associative
\(R\)-algebra, not necessarily commutative. For \(a,b\in A\), put
\[
  [a,b]=ab-ba,
  \qquad
  a\cdot b=L_a(b)=R_b(a).
\]
Then
\[
\begin{aligned}
  (\operatorname{ad}x^p)(y)
  &=[x^p,y] \\
  &=(L_x^p-R_x^p)(y) \\
  &=(L_x-R_x)^p(y) \\
  &=(\operatorname{ad}x)^p(y),
\end{aligned}
\]
which gives Jacobson's first formula:
\[
  \tag{i}\label{sga3:vii:paragraph:5-1:equation:i}
  \operatorname{ad}(x^p)
  =
  (\operatorname{ad}x)^p.
\]

If \(a_1,\ldots,a_p\) are arbitrary elements of \(A\), and \(N\)
is the symmetrization operator (cf.\ \SGARefLink{sga3:VIIA:4.2}{4.2}), then
\begin{equation}
  \tag{*}\label{sga3:vii:paragraph:5-1:equation:star}
\begin{aligned}
  N(a_1\otimes\cdots\otimes a_p)
  &=
  \sum_{\sigma\in\mathfrak S_p}
    a_{\sigma(1)}\cdots a_{\sigma(p)} \\
  &=
  \sum_{\tau\in\mathfrak S_{p-1}}
  [a_{\tau(1)},
    [a_{\tau(2)},
      [\cdots,
        [a_{\tau(p-1)},a_p]
      \cdots]
    ]
  ].
\end{aligned}
\end{equation}
Here \(\sigma\) ranges over the permutations of \(p\) letters and
\(\tau\) over those of \(p-1\) letters.

Indeed, the final expression is equal to
\[
  \sum_{\tau}
  \sum_{r=0}^{p-1}
  \sum_{i_1<\cdots<i_r}
  (-1)^s
  a_{\tau(i_1)}\cdots a_{\tau(i_r)}
  a_p
  a_{\tau(j_s)}\cdots a_{\tau(j_1)},
\]
where \(\tau\) ranges over the permutations of \(p-1\) letters;
\(i_1,\ldots,i_r\) ranges over the strictly increasing sequences
in \([1,p-1]\); and \(j_1,\ldots,j_s\) is the strictly increasing
sequence of those integers of \([1,p-1]\) which are different
from \(i_1,\ldots,i_r\). For fixed \(r\), the sum of the terms
\[
  (-1)^s
  a_{\tau(i_1)}\cdots a_{\tau(i_r)}
  a_p
  a_{\tau(j_s)}\cdots a_{\tau(j_1)}
\]
is plainly
\[
  (-1)^s\binom{p-1}{s}
  \sum_{\rho}
  a_{\rho(1)}\cdots a_{\rho(r)}
  a_p
  a_{\rho(r+1)}\cdots a_{\rho(p-1)},
\]
where \(\rho\) ranges over the permutations of \(p-1\) letters.
But
\[
  (-1)^s\binom{p-1}{s}=1
  \quad\text{in }\mathbb F_p,
\]
because in \(\mathbb F_p[x]\),
\[
  (x-1)^p=x^p-1=(x-1)(x^{p-1}+\cdots+1),
\]
and hence
\[
  (x-1)^{p-1}=x^{p-1}+\cdots+1.
\]
This proves \((*)\).

On the other hand, if \(x_0,x_1\in A\), then
\[
  (x_0+x_1)^p
  =
  x_0^p+x_1^p+
  \sum_z x_{z(1)}x_{z(2)}\cdots x_{z(p)},
\]
where \(z\) ranges over the nonconstant maps
\([1,p]\to\{0,1\}\). It follows that
\[
\begin{aligned}
  (x_0+x_1)^p
  ={}&
  x_0^p+x_1^p \\
  &+
  \sum_{0<r<p}
  \frac{1}{r!(p-r)!}\,
  N(
    \underbrace{x_0,\ldots,x_0}_{r},
    \underbrace{x_1,\ldots,x_1}_{p-r}
  ).
\end{aligned}
\]
\footnote{The editors inserted the following explanation from
\textup{[DG70]}, III.7, \SGARefLink{sga3:VIIB:3.2:prop}{3.2}.}
By \((*)\),
\[
\begin{aligned}
  &N(
    \underbrace{x_0,\ldots,x_0}_{r},
    \underbrace{x_1,\ldots,x_1}_{p-r}
  ) \\
  &\qquad =
  r!(p-1-r)!
  \sum_t
  [x_{t(1)},
    [x_{t(2)},
      [\cdots,
        [x_{t(p-1)},x_1]
      \cdots]
    ]
  ],
\end{aligned}
\]
where \(t\) ranges over the maps
\([1,p-1]\to\{0,1\}\) which take the value \(0\) exactly \(r\)
times. We obtain Jacobson's second formula:
\begin{equation}
  \tag{ii}\label{sga3:vii:paragraph:5-1:equation:ii}
\begin{aligned}
  (x_0+x_1)^p
  ={}&
  x_0^p+x_1^p \\
  &-
  \sum_{0<r<p}\sum_t\frac1r
  [x_{t(1)},
    [x_{t(2)},
      [\cdots,
        [x_{t(p-1)},x_1]
      \cdots]
    ]
  ],
\end{aligned}
\end{equation}
where \(t\) has the same meaning.

\medskip
\noindent\textbf{5.2.}
\phantomsection
\label{sga3:VIIA:5.2}

Let now \(\mathfrak g\) be an \(R\)-Lie algebra. A map
\[
  x\longmapsto x^{(p)}
\]
from \(\mathfrak g\) to itself is said to make \(\mathfrak g\) a
\emph{restricted \(p\)-Lie algebra over \(R\)} if the following
conditions hold:
\[
\begin{array}{ll}
  (0) &
  (\lambda x)^{(p)}
  =
  \lambda^p x^{(p)}
  \quad
  (\lambda\in R,\ x\in\mathfrak g), \\[0.4em]
  (i) &
  \operatorname{ad}x^{(p)}
  =
  (\operatorname{ad}x)^p
  \quad
  (x\in\mathfrak g),
\end{array}
\]
and
\[
\begin{aligned}
  (ii)\quad
  (x_0+x_1)^{(p)}
  ={}&
  x_0^{(p)}+x_1^{(p)} \\
  &-
  \sum_{0<r<p}\sum_t\frac1r
  [x_{t(1)},
    [x_{t(2)},
      [\cdots,
        [x_{t(p-1)},x_1]
      \cdots]
    ]
  ],
\end{aligned}
\]
where \(t\) ranges over the maps
\([1,p-1]\to\{0,1\}\) which take the value \(0\) exactly \(r\)
times, and \(x_0,x_1\in\mathfrak g\). The map
\(x\mapsto x^{(p)}\) is called the \emph{symbolic \(p\)-th power}.

For example, if \(A\) is an associative \(R\)-algebra, \SGARefLink{sga3:VIIB:5.1:proposition}{5.1} shows
that the \(R\)-module underlying \(A\), endowed with
\[
  [x,y]=xy-yx,
  \qquad
  x^{(p)}=x^p,
\]
is a \(p\)-Lie algebra, denoted \(A_{\mathrm{Lie}}\). We call it
the \(p\)-Lie algebra underlying \(A\).

We shall chiefly consider restricted \(p\)-Lie subalgebras of
algebras of the form \(A_{\mathrm{Lie}}\). Here is one example.
Let \(S\) be a scheme of characteristic \(p>0\), and let \(X\) be
an \(S\)-scheme. Recall that a derivation of \(X\) over \(S\) is
an endomorphism \(D\) of the sheaf of abelian groups
\(\mathcal O_X\) such that
\[
  D(\lambda s)=\lambda D(s),
  \qquad
  D(st)=(Ds)t+s(Dt),
\]
whenever \(\lambda\), respectively \(s,t\), are sections of
\(\mathcal O_S\), respectively \(\mathcal O_X\), on open subsets
for which the formulas make sense. Leibniz's formula
\[
  D^n(st)
  =
  \sum_{i=0}^n
  \binom ni
  (D^is)(D^{n-i}t)
\]
shows that \(D^p\) is again a derivation of \(X\) over \(S\), since
\[
  \binom pi\equiv0\pmod p
  \qquad(i\ne0,p).
\]
It follows that:

\begin{quote}
\textit{The algebra \(\operatorname{Der}_{X/S}\) of derivations of
\(X\) over \(S\) is a restricted \(p\)-Lie subalgebra of the
\(\Gamma(S,\mathcal O_S)\)-algebra of differential operators of
\(X\) over \(S\).}
\end{quote}

\medskip
\noindent\textbf{5.2.1.}
\phantomsection
\label{sga3:VIIA:5.2.1}
If \(\mathfrak g\) and \(\mathfrak h\) are two restricted
\(p\)-Lie algebras, a \emph{homomorphism}
\[
  h:\mathfrak g\longrightarrow\mathfrak h
\]
is an \(R\)-linear map such that, for \(x,y\in\mathfrak g\),
\[
  h([x,y])=[h(x),h(y)],
  \qquad
  h(x^{(p)})=h(x)^{(p)}.
\]
The composite of two homomorphisms is again a homomorphism, so we
may speak of the \emph{category of restricted \(p\)-Lie algebras
over \(R\)}.

If \((X,\mathcal R)\) is a ringed space, we say that an
\(\mathcal R\)-module \(\mathfrak g\) is endowed with the
structure of a restricted \(p\)-Lie algebra over \(\mathcal R\) if,
for every open subset \(U\), the module
\(\Gamma(U,\mathfrak g)\) is a restricted \(p\)-Lie algebra over
\(\Gamma(U,\mathcal R)\), and if all restriction maps are
homomorphisms.

